% !TeX document-id = {9e81b2f4-72a3-4a11-b845-f09d8e123456}
% !TeX program = pdflatex
% !TeX TXS-program:compile = txs:///pdflatex/[-shell-escape]
% !TeX TXS-program:pdflatex = pdflatex -synctex=1 -interaction=nonstopmode --shell-escape %.tex

\documentclass[../../main_compilation.tex]{subfiles}

\begin{document}

\chapter{Knowledge Base \& Bitácora de Aprendizaje}
\label{chap:knowledge_base}
\setCustomHeader{{\color{csecPurple}\faBrain}\ \textbf{Playbook}}{Knowledge Base \& Bitácora}{\today}

Este capítulo constituye el repositorio central de conocimientos técnicos, conceptos teóricos y metodologías asimiladas durante auditorías, resolución de laboratorios y estudio bibliográfico. A diferencia de los reportes individuales, aquí la información se estructura de forma \textbf{temática y transversal}, vinculando cada habilidad con los libros, sitios web, documentación y máquinas específicas donde fue adquirida.

% ================================================================================
% 1. Metodología de Registro y Fuentes de Estudio
% ================================================================================
\section{Metodología de Registro y Fuentes de Consulta}
Para que el conocimiento perdure y sea consultable de manera inmediata en auditorías reales o exámenes prácticos (como OSCP o CPTS), cada concepto registrado debe responder a tres preguntas fundamentales:
\begin{enumerate}
	\item \textbf{¿Qué aprendí?} (Concepto técnico, mecanismo de funcionamiento, payload o comando clave).
	\item \textbf{¿Dónde lo aprendí?} (Libro con autor y capítulo, sitio web/documentación con URL, o máquina vulnerable con enlace interno).
	\item \textbf{¿Cómo se aplica?} (Comando reproducible, script o prueba de concepto).
\end{enumerate}

\begin{infobox}[Badges de Fuentes Disponibles]
	Para estandarizar el origen del conocimiento, dispones de la macro \texttt{\textbackslash sourceBadge\{\ldots\}} y el comando \texttt{\textbackslash resourceItem\{Tipo\}\{Título/Autor\}\{Detalle/Enlace\}}:
	\begin{center}
		\small
		\sourceBadge{Libro} \hspace{4pt}
		\sourceBadge{Web} \hspace{4pt}
		\sourceBadge{Docs} \hspace{4pt}
		\sourceBadge{Curso} \hspace{4pt}
		\sourceBadge{Lab} \hspace{4pt}
		\sourceBadge{Writeup} \hspace{4pt}
		\sourceBadge{Video} \hspace{4pt}
		\sourceBadge{Paper}
	\end{center}
\end{infobox}

% ================================================================================
% 2. Fundamentos de Lenguajes & Arquitectura
% ================================================================================
\section{Fundamentos de Lenguajes \& Arquitectura}

\subsection{Assembly (ASM x86\_64) \& Arquitectura de Computadores}
El conocimiento de ensamblador x86\_64 es fundamental para tareas de ingeniería inversa, depuración de binarios (\textit{pwn}) y desarrollo de exploits.

\begin{learningbox}[Conceptos Clave Aprendidos: Arquitectura x86\_64]
	\begin{itemize}[leftmargin=*]
		\item \textbf{Registros de Propósito General:} Familiarización con los registros de 64 bits (\texttt{rax}, \texttt{rbx}, \texttt{rcx}, \texttt{rdx}, \texttt{rsi}, \texttt{rdi}, \texttt{rsp}, \texttt{rbp}, \texttt{r8}--\texttt{r15}) y sus sub-registros de 32, 16 y 8 bits (\texttt{eax}, \texttt{ax}, \texttt{al}).
		\item \textbf{Convenio de Llamadas (System V AMD64 ABI):} Los primeros seis argumentos enteros/punteros a funciones en Linux se pasan a través de registros en orden estricto: \texttt{rdi}, \texttt{rsi}, \texttt{rdx}, \texttt{rcx}, \texttt{r8}, \texttt{r9}. Los argumentos adicionales van a la pila (\textit{stack}).
		\item \textbf{Llamadas al Sistema (\texttt{syscall}):} El número de syscall se coloca en \texttt{rax}, los argumentos en \texttt{rdi}, \texttt{rsi}, etc., y el valor de retorno se recupera en \texttt{rax}.
		\item \textbf{Estructura del Stack Frame:} Creación de prólogo (\texttt{push rbp; mov rbp, rsp}) y epílogo (\texttt{leave; ret}) y reserva de variables locales restando a \texttt{rsp}.
	\end{itemize}
\end{learningbox}

\begin{sourcebox}[Fuentes y Bibliografía Utilizadas]
	\begin{itemize}[leftmargin=*]
		\resourceItem{Libro}{Practical Binary Analysis (Dennis Andriesse)}{Capítulo 2: \textit{x86 and x86-64 Architecture Overview} (No Starch Press)}
		\resourceItem{Libro}{Assembly Language Step-by-Step (Jeff Duntemann)}{Fundamentos de registros, banderas de estado y direccionamiento de memoria}
		\resourceItem{Docs}{x86-64 Cheat Sheet — Brown University}{\url{https://cs.brown.edu/courses/cs033/docs/guides/x64_cheatsheet.pdf}}
		\resourceItem{Web}{Linux System Call Table for x86\_64 (Ray Chapman)}{\url{https://blog.rchapman.org/posts/Linux_System_Call_Table_for_x86_64/}}
		\resourceItem{Docs}{Intel 64 and IA-32 Architectures Software Developer's Manual}{\url{https://www.intel.com/content/www/us/en/developer/articles/technical/intel-sdm.html}}
	\end{itemize}
\end{sourcebox}

\begin{codebox}[nasm]{Ejemplo Práctico: Syscall sys\_write en x86\_64 (hello.asm)}
	global _start

	section .text
	_start:
	; syscall: sys_write (rax = 1, rdi = fd 1, rsi = buffer, rdx = len)
	mov rax, 1                  ; ID de syscall sys_write en Linux x86_64
	mov rdi, 1                  ; File Descriptor: 1 (stdout)
	mov rsi, msg                ; Puntero al mensaje
	mov rdx, len                ; Longitud en bytes
	syscall

	; syscall: sys_exit (rax = 60, rdi = status code 0)
	mov rax, 60
	xor rdi, rdi
	syscall

	section .data
	msg db "Hello, x86_64 Assembly!", 0x0A
	len equ $ - msg
\end{codebox}

\subsection{Bash Scripting \& Automatización de Tareas}
Bash es el lenguaje nativo de los entornos Linux para automatización rápida de auditorías, parsing de ficheros de log y creación de herramientas de reconocimiento al vuelo.

\begin{learningbox}[Conceptos Clave Aprendidos: Bash Avanzado]
	\begin{itemize}[leftmargin=*]
		\item \textbf{Redirección y Descriptores de Ficheros:} Control preciso de \texttt{stdin} (0), \texttt{stdout} (1) y \texttt{stderr} (2). Uso de \texttt{2>\&1} para unificar flujos y \texttt{/dev/null} para silenciar ruido en bucles.
		\item \textbf{One-Liners de Reconocimiento:} Uso de combinaciones de \texttt{curl}, \texttt{grep}, \texttt{sed}, \texttt{awk} y \texttt{xargs} para extraer subdominios o endpoints activos en segundos sin herramientas pesadas.
		\item \textbf{Subshells y Captura de Salida:} Diferencia entre asignación síncrona mediante \texttt{\$(comando)} y ejecución en subshell background con \texttt{\&} junto al comando \texttt{wait}.
		\item \textbf{Manejo de Señales (\texttt{trap}):} Limpieza automática de ficheros temporales y sockets al pulsar \texttt{Ctrl+C} (\texttt{SIGINT}) en scripts interactivos.
	\end{itemize}
\end{learningbox}

\begin{sourcebox}[Fuentes y Bibliografía Utilizadas]
	\begin{itemize}[leftmargin=*]
		\resourceItem{Libro}{The Linux Command Line: A Complete Introduction (William Shotts)}{Capítulos 20 al 35: Entornos shell, variables, control de flujo y scripts robustos}
		\resourceItem{Libro}{Wicked Cool Shell Scripts (Dave Taylor \& Brandon Perry)}{2ª Edición — Scripts prácticos para administración de sistemas y seguridad}
		\resourceItem{Web}{Bash Hackers Wiki — Guía profunda de sintaxis y buenas prácticas}{\url{https://wiki.bash-hackers.org}}
		\resourceItem{Web}{ExplainShell — Análisis visual e interactivo de comandos y flags}{\url{https://explainshell.com}}
		\resourceItem{Docs}{GNU Bash Reference Manual Oficial}{\url{https://www.gnu.org/software/bash/manual/}}
	\end{itemize}
\end{sourcebox}

\begin{bashbox}{Snippet Útil: Escaneo Rápido de Hosts Vía Ping en Bash}
	#!/usr/bin/env bash
	# Sweep rápido de red /24 en paralelo usando subshells
	SUBNET="10.10.10"
	echo "[*] Escaneando subred ${SUBNET}.0/24..."

		for ip in $(seq 1 254); do
	(ping -c 1 -W 1 "${SUBNET}.${ip}" &>/dev/null && echo "[+] Host activo: ${SUBNET}.${ip}") &
	done
	wait
	echo "[*] Barrido completado."
\end{bashbox}

% ================================================================================
% 3. Técnicas de Explotación & Casos Prácticos
% ================================================================================
\section{Técnicas de Explotación \& Seguridad Ofensiva}
En esta sección se correlacionan las técnicas de intrusión con los retos o laboratorios resueltos dentro de este mismo documento.

\subsection{Explotación de Servicios de Red: Samba Username Map Script}
\begin{learningbox}[Mecanismo de Vulnerabilidad (CVE-2007-2447)]
	En versiones vulnerables de Samba ($3.0.0$ a $3.0.25\text{rc}3$), la directiva \texttt{username map script} ejecutaba scripts de shell externos concatenando directamente el nombre de usuario sin escapar metacaracteres (\texttt{`} o \texttt{\$(...)}). Al pasar un payload como nombre de usuario vía SMB, el comando se ejecuta con privilegios de \texttt{root}.
\end{learningbox}

\begin{sourcebox}[Fuentes y Referencias del Concepto]
	\begin{itemize}[leftmargin=*]
		\resourceItem{Writeup}{\hyperref[sec:htb-lame]{HackTheBox — Lame (Capítulo~\ref{chap:hackthebox})}}{Máquina donde se identificó y explotó manualmente este vector en la práctica}
		\resourceItem{Docs}{NIST NVD — Registro Oficial CVE-2007-2447}{\url{https://nvd.nist.gov/vuln/detail/CVE-2007-2447}}
		\resourceItem{Web}{Rapid7 Metasploit Module Documentation}{\url{https://www.rapid7.com/db/modules/exploit/multi/samba/usermap_script/}}
	\end{itemize}
\end{sourcebox}

\subsection{Explotación de Binarios: Format String Vulnerability}
\begin{learningbox}[Fuga Arbitraria de Memoria del Stack]
	Cuando una función variádica como \texttt{printf()} recibe una entrada del usuario sin una cadena de formato estricta (\texttt{printf(buf)} en lugar de \texttt{printf("\%s", buf)}), el usuario puede suministrar especificadores como \texttt{\%p} o \texttt{\%x} para leer valores del stack, o \texttt{\%n} para escribir en memoria. En x86\_64 es fundamental considerar el orden Little-Endian al reconstruir cadenas desde valores hexadecimales.
\end{learningbox}

\begin{sourcebox}[Fuentes y Referencias del Concepto]
	\begin{itemize}[leftmargin=*]
		\resourceItem{Lab}{\hyperref[sec:ctf-pico-format-string-1]{picoCTF — Format String 1 (Capítulo~\ref{chap:ctfs})}}{Reto donde se aplicó extracción de datos posicionales \texttt{\%14\$p} para leer la bandera secreta}
		\resourceItem{Libro}{Hacking: The Art of Exploitation (Jon Erickson)}{Capítulo 0x300: \textit{Exploitation — Format Strings}}
		\resourceItem{Web}{OWASP Format String Vulnerability Reference}{\url{https://owasp.org/www-community/attacks/Format_string_attack}}
	\end{itemize}
\end{sourcebox}

% ================================================================================
% 4. Matriz Resumen de Conocimientos Adquiridos
% ================================================================================
\section{Matriz Global de Aprendizajes \& Fuentes}
Esta tabla proporciona un índice rápido para localizar de inmediato qué técnica domina el auditor y dónde consultar la referencia teórica o práctica correspondiente:

\begin{table}[h!]
	\centering
	\small
	\begin{tabularx}{\textwidth}{p{4.0cm} p{2.8cm} X c}
		\toprule
		\rowcolor{csecDark}
		\textbf{\color{white}Técnica / Concepto}                                                                                  &
		\textbf{\color{white}Categoría}                                                                                           &
		\textbf{\color{white}Fuente Principal (Libro / Web / Lab)}                                                                &
		\textbf{\color{white}Referencia}                                                                                            \\
		\midrule
		\textbf{Llamadas al sistema y registros en x86\_64}                                                                       &
		Arquitectura / ASM                                                                                                        &
		\sourceBadge{Libro} \textit{Practical Binary Analysis} \newline \sourceBadge{Web} Brown Univ. Cheatsheet                  &
		\hyperref[chap:knowledge_base]{Cap.~\ref{chap:knowledge_base}}                                                              \\
		\midrule
		\textbf{Subshells y Descriptores de Ficheros}                                                                             &
		Scripting / Bash                                                                                                          &
		\sourceBadge{Libro} \textit{The Linux Command Line} \newline \sourceBadge{Docs} GNU Bash Manual                           &
		\hyperref[chap:knowledge_base]{Cap.~\ref{chap:knowledge_base}}                                                              \\
		\midrule
		\textbf{Inyección de comandos vía Samba}                                                                                  &
		Redes / RCE                                                                                                               &
		\sourceBadge{Writeup} Máquina \textbf{Lame} (HTB) \newline \sourceBadge{Docs} CVE-2007-2447                               &
		\hyperref[sec:htb-lame]{Sec.~\ref{sec:htb-lame}}                                                                            \\
		\midrule
		\textbf{Format String Stack Leak (\%p)}                                                                                   &
		Binary Exploitation                                                                                                       &
		\sourceBadge{Lab} Reto \textbf{Format String 1} (picoCTF) \newline \sourceBadge{Libro} \textit{Hacking: The Art of Expl.} &
		\hyperref[sec:ctf-pico-format-string-1]{Sec.~\ref{sec:ctf-pico-format-string-1}}                                            \\
		\bottomrule
	\end{tabularx}
	\caption{Matriz consolidada de conocimientos adquiridos y fuentes de procedencia.}
\end{table}

\end{document}
